%!TEX root = ../../main.tex
%% 1.6 논문의 구성
\section{논문의 구성}
\label{sec:intro-outline}

제 \ref{ch:related} 장은 관련 연구와 본 논문이 쓰는 개념의 정의·기원을 정리하고 선행 연구의 위치를 밝힌다. 제 \ref{ch:data} 장은 비동기 공개 텔레메트리의 구조, 원천 코퍼스, 데이터 역할, 교전 전 입력의 구성을 다룬다(M-RQ1). 제 \ref{ch:engagement} 장은 교전 사례의 정의, 상수의 출처, 규모 계급과 코호트, 결과 시점 규칙, 정의가 볼 수 없는 것을 다룬다(M-RQ1). 제 \ref{ch:value} 장은 동결 평가기 $\Vhat$와 $\SVI$의 정의, 그리고 그 측정이 무엇을 재는지에 대한 검증을 제시한다(M-RQ2). 제 \ref{ch:prediction} 장은 예측기 $q$, 기준선, 선정·보정 절차와 코호트별 주 결과, 전이·합동 대비를 제시한다(M-RQ3). 제 \ref{ch:validation} 장은 균형 구간, 작은 변화 제외, 점수 분해, 외부 전이로 예측의 의미와 범위를 검증한다(M-RQ4). 제 \ref{ch:discussion} 장은 네 단계가 어떻게 연결되는지, 선행 연구 대비 기여, 금지 주장, 한계, 향후 과제를 통합하여 논의하고 결론을 맺는다. 부록에는 재현 자료, 상수표, 데이터 역할 census, 수치의 증거 추적표, 이연 계획을 둔다.
